\hypertarget{group-state-and-general-measurement}{%
\section{Group state and general measurement}\label{method-insurance-methods-group-state-and-general-measurement}}

The following cases apply selected general-measurement and service-allocation provisions from the retained AASB 17 compilation \citep{aasb17}. They extend the premium-allocation example without presuming actuarial calibration or a complete insurer model.

The general method distinguishes fulfilment cash flows (FCF), including an
explicit risk adjustment (RA), from the contractual service margin (CSM),
which holds unearned profit. The liability for remaining coverage (LRC)
concerns future service; the liability for incurred claims (LIC) concerns
service and insured events already experienced. Keeping these components
separate makes a change in future estimates different from an unpaid claim.

Contract records retain insured risks, premiums/benefits, service and settlement dates, options/guarantees, cancellation/repricing rights and acquisition costs. Group admission evaluates portfolio/similar-risk management, onerous status and required issuance cohorts under AASB 17.14--24, recognition under 25--28 and boundary under 34. Cash outside the boundary belongs to future contracts, not this group's profit. GroupState retains expected dated cash flows, current discount curves, inception rates, explicit risk adjustment, CSM, loss component, LRC/LIC, coverage units, service experience and assumption vintages. Same premium totals are not sufficient state for merging different guarantees or claim tails.

General measurement follows 32--38: fulfilment cash flows are probability-weighted future net outflows, adjusted for timing/financial risk without double counting, plus explicit non-financial risk adjustment. CSM removes initial unearned gain under its conditions; an initial net loss is recognized under 47, not negative CSM. Paragraphs 40--44 separate remaining coverage and incurred claims. For a nonparticipating group, accrete CSM at the specified locked-in rate, apply future service changes at their required rate to CSM up to its available amount, record any excess loss and allocate current service over current plus future coverage units under B119. B96--B97 distinguish future service from incurred-claim changes, financial changes and current experience; not every revision is a CSM adjustment.

A MeasurementInput must resolve cash-flow/law, discount, risk-adjustment method, service grouping and coverage-unit evidence. Missing probabilities or RA support permit conditional sensitivity cases but block a source-supported measurement. Risk adjustment is the compensation the entity requires under 37; a standard deviation, confidence quantile or arbitrary percentage is only a candidate method until justified for that entity. Disclosures and the confidence-level equivalent, where applicable, are separate obligations. Market-consistent financial inputs and entity risk compensation must not be conflated.

\hypertarget{complete-two-year-general-measurement-case}{%
\section{Complete two-year general-measurement case}\label{method-insurance-methods-complete-two-year-general-measurement-case}}

Take a two-year nonparticipating group with one premium 120 received at inception, claims expected/actually incurred and paid 40 each year, zero discount rate, RA 10 initially released 5 with each equal service year, no acquisition costs, investment component, lapses, taxes, finance changes or reinsurance. Coverage units are 1 in each year because equal benefits and coverage are assumed. Before premium payment, expected net cash outflow is 80\ensuremath{-}120=\ensuremath{-}40, RA 10, FCF \ensuremath{-}30 and CSM 30, with net initial carrying amount zero. After the premium debit cash 120/ credit LRC 120. Its components are remaining expected claims 80, RA 10 and CSM 30.

Each year releases expected service claims 40, RA 5 and CSM 15: debit LRC 60/ credit insurance revenue 60. Incur claims 40: debit service expense 40/credit LIC 40; payment debits LIC 40/credits cash 40. End of year one: cash 80, LRC 60 (claims 40, RA 5, CSM 15), LIC 0 and equity/profit 20. End year two: cash 40, liabilities 0, cumulative equity/profit 40. Revenue 120 and claims 80 reconcile to premium 120 less claims cash 80. No investment component has been mistaken for revenue.

At year-one close, a separate pre-next-service change increases the remaining claim estimate from 40 to 60, RA unchanged. FCF rises 20; consume remaining CSM 15 and recognize service loss 5. LRC rises from 60 to 65, comprising FCF 65, CSM 0 and loss component 5 as a tracked portion, not an added liability on top. Debit service loss 5/credit LRC 5 with component transfer CSM\ensuremath{\rightarrow}FCF 15. A decrease of 8 in the revised future claim estimate, before any additional service/finance movement, first reverses loss 5 and restores CSM 3. FCF 57 plus CSM 3 gives LRC 60. Credit loss reversal 5/debit LRC 5. This tests 48--52 ordering. Later systematic allocation of loss component to releases/finance must reduce it to zero by coverage end; it cannot be amortized as revenue a second time.

In a separate case, an already-incurred unpaid claim estimate rises 40 to 48. Debit service expense 8/credit LIC 8; do not consume future-service CSM. A discount-rate change is financial under B97; its profit/OCI route follows the selected permitted finance policy. The zero-rate base case does not remove the need for separate current and locked-in rate calculations in a real group.

For the unreversed onerous alternative, carry the loss component through the remaining service year. Opening LRC 65 comprises future claims 60 and RA 5, with CSM 0 and loss component 5 within that LRC. A declared proportional allocation assigns 5/65 of each expected claim/RA release to the loss component, consistent with paragraphs 49--52 and B124's revenue exclusions. The expected claim release allocates 60/13 to loss reversal and the RA release 5/13; the total is 5, leaving loss component 0. Debit LRC 65, credit revenue 60 and credit loss reversal 5. Separately incur and pay claims 60 through LIC. Period profit is 60\ensuremath{-}60+5=5; combined with first-year profit 20 and the earlier onerous loss 5, lifetime profit is 20. Closing cash 20 equals equity 20 and all liabilities are zero. Treating all 65 as revenue while also reversing loss 5 overstates profit by 5.

\hypertarget{nonzero-discount-and-service-allocation}{%
\section{Nonzero discount and service allocation}\label{method-insurance-methods-nonzero-discount-and-service-allocation}}

A second complete case keeps premium 120, two equal coverage units and RA 10, but expected and realized claims are 44 after one year and 48.4 after two years. Both the inception and current discount curves remain a flat 10\% throughout; there are no new estimates, options, fees, taxes, investments or OCI election. This RA schedule is an authored measurement input, not an actuarial calibration: it declines by 5 per year, with its entire change in the service result under 81. All other insurance finance effects go to profit or loss. The retained source supports the mechanics in 44, 80--88, B119 and B123--B124.

At inception the present value of the two expected claims is
\[
 \frac{44}{1.1}+\frac{48.4}{1.1^2}=40+40=80.
\]
 FCF before premium is 80+10\ensuremath{-}120=\ensuremath{-}30, CSM 30, and after premium LRC 120. In year one the claim PV accretes 8; CSM accretes 3 at the locked-in rate. Allocate the resulting CSM 33 across the two current/future units, releasing 16.5. Revenue is expected claims 44 plus RA 5 plus CSM 16.5=65.5. Service expense 44 and finance expense 11 leave profit 10.5. After paying 44, cash 76 and LRC 65.5 (claim PV 44, RA 5, CSM 16.5) reconcile to equity 10.5.

In year two claim PV accretes 4.4 and CSM accretes 1.65. Release all remaining CSM 18.15, RA 5 and expected claims 48.4, giving revenue 71.55. Service expense 48.4 and finance expense 6.05 leave profit 17.1. Final cash/equity 27.6 equals premium 120 minus claims 92.4; all group liabilities are zero. Total revenue 137.05 less finance expense 17.05 equals premium 120. Setting revenue equal to cash premium each year or releasing 15 of CSM each year would ignore the specified financing. These are separate component calculations followed by debit finance expense/ credit LRC for accretion, debit LRC/credit revenue for service and the stated LIC recognition/settlement entries. Cash is deliberately assumed uninvested.

\hypertarget{paa-reinsurance-and-valuation-methods}{%
\section{PAA, reinsurance and valuation methods}\label{method-insurance-methods-paa-reinsurance-and-valuation-methods}}

The preceding premium-allocation case states separate eligibility and recognition timing for gross insurance and reinsurance. PAA under 53--59 requires the one-year coverage or nonmaterial-difference condition; a short payment lag is not the eligibility test. LRC starts with received premium and applicable acquisition adjustments, then releases insurance revenue under the prescribed service pattern. LIC still measures incurred obligations using the applicable fulfilment-cash-flow rules; the allowed short-settlement discount simplification needs its own conditions. Emerging onerous facts require a loss assessment; a premium balance cannot be presumed sufficient.

A one-year policy can leave claims unpaid for several years. The distinction
between the period of insurance service and the period until settlement is
therefore essential even when PAA is available. The official ASBE~25 text
separates three timing tests. Article~50 permits PAA when each contract's
coverage is no longer than one year, or when the stated no-material-difference
condition holds. Article~54 concerns a different interval: the time between
each portion of service and the related premium due date. If that interval
does not exceed a year, the remaining-coverage liability need not be adjusted
for financing. Article~56 instead measures from claim incurrence to expected
payment or receipt. Its option to omit the time-value and financial-risk
adjustment requires settlement within a year of incurrence
\citep[Articles 50, 54, 56]{mofasbe25}. The retained AASB~17 compilation makes
the corresponding distinctions in paragraphs 53, 56 and 59(b).

Consider a separate group of issued policies with exactly one year of coverage.
All amounts are illustrative currency units. Premium 120 is due and received
when coverage starts, with no earlier onerous trigger; group recognition is
therefore at that date under Article~16. Assume valid insurance grouping,
uniform risk release, no acquisition costs, investment component, reinsurance,
tax, other expenses or other movements, and no facts indicating onerous
remaining coverage. Elect the permitted remaining-coverage financing exception.
After half a year, recognise revenue 60 and retain remaining coverage of 60.
Articles~53 and 57 support this roll-forward and time allocation; Article~57
requires a different allocation when risk release is not predominantly
time-based \citep{mofasbe25}. These assumptions isolate the claim measurement
from other changes in the group's result.

At that reporting instant a claim is first incurred, with an expected payment
of 60.5 exactly two years later. This is a mean estimate, not a certain
settlement. Assume an appropriate flat annual effective discount rate of
10\% and a separately estimated non-financial risk adjustment of 5 in current
measurement units. The inputs are authored for the example; applying them to
an insurer would require market and actuarial evidence. Articles~21, 23, 25
and 26 require current probability-weighted cash flows, appropriate discounting
and a separate non-financial risk adjustment \citep{mofasbe25}. The claim's
cash-flow present value is $60.5/1.1^2=50$, so its liability is $50+5=55$.
The risk adjustment is not discounted a second time. The claim has only just
been incurred, so no intervening finance accrual enters this measurement.

Record debit cash 120/credit remaining coverage 120 on premium receipt, then
debit remaining coverage 60/credit insurance revenue 60 for the service
provided. Recognising the claim gives debit insurance service expense 55/credit
incurred claims 55, consistent with Article~31. The group then has cash 120,
remaining coverage 60, incurred claims 55 and a cumulative result of 5:
$120-60-55=60-55=5$. These are balances attributable to the group at mid-year,
not a complete insurer forecast. Using the undiscounted $60.5+5=65.5$ for this
two-year claim would overstate the required liability by 10.5 and understate
the result by the same amount. The shortcut is wrong under the stated rules
and assumptions. A zero rate would remove that discounting difference while
leaving the separate risk adjustment intact. Articles~86(3) and 97 also require
disclosure of the distinct claim components and the PAA/financing choices
\citep{mofasbe25}.

This comparison uses identified editions. The Ministry page reproduces the
2020 ASBE~25 text; Article~120 states 1 January 2023, but that sentence does
not settle every entity group's adoption schedule.
Section~\ref{sec:insurance-adoption} distinguishes the Ministry's later
explanation of that timetable from a company's reported adoption. The retained
December 2022 AASB~17 compilation applies to annual periods beginning on or
after 1 January 2023 but before 1 July 2026. The shared result in this example
establishes a scoped measurement comparison. Applying it to a particular entity
still requires its applicable rules and contract facts; later amendments and
broader equivalence require separate evidence.

Reinsurance has its own group, boundary, recognition time, counterparty risk, cash-flow, CSM/cost and possibly loss-recovery component under AASB~17 paragraphs 60--70A. Do not multiply every gross insurance liability by a quota share: excluded claims, limits, deductibles, timing, reinstatements and counterparty non-performance matter. Claim 100 with an eligible 50\% recovery but expected 90\% collectability gives conditional expected recovery 45, not 50; altering transferred insurance risk and reinsurer non-performance are different inputs. This simple calculation does not establish treaty recognition or RA by itself. Gross liability 100 remains; reinsurance recovery is a separate asset/measurement with its own authority.

For direct participation, test B101's conditions using the actual underlying items, substantial-share and variation judgments. Retain underlying fair values, entity fee share and guarantee/risk-mitigation choices; paragraphs 45 and B101--B118 specify the variable-fee adjustments to CSM, with onerous and mitigation exceptions. A local component illustration of a qualifying share: underlying-value increase 20 and entity's fixed 10\% share gives candidate variable-fee change 2; eligibility, guarantees and other FCF changes must be evaluated before posting. The full option/guarantee valuation and transition calculations are further supported profiles, with source and numerical obligations, not an assumed linear fee model.

Eligibility depends on the contract and the insurer's expectations at inception.
The policyholder must participate in a clearly identified pool, with the
required substantial share of fair-value returns and substantial variability
in payments. The variability assessment covers the contract's duration on a
probability-weighted present-value basis. A minimum guarantee can bind in some
scenarios without automatically excluding direct participation. Later market
movements do not reopen the classification; contract modification follows its
own rules. The identified reference pool also need not be held by the insurer,
so the contractual link and the actual asset holdings require separate records
\citep[paragraphs B101--B108 and 72]{aasb17}.

\clearpage
A small settlement example shows why the guarantee matters even when the
stated share is fixed. Consider only a maturity-benefit component, in one
currency at one date. Assume a held pool allocated to this component has amount
$U$, and the policyholder is owed the larger of 90 and $0.9U$. These are
illustrative contractual choices. The amount left from the pool after payment
is $U-\max(90,0.9U)$. Writing the payment as
$0.9U+\max(90-0.9U,0)$ shows that the amount left is
$0.1U-\max(90-0.9U,0)$: the nominal retained share less the extra payment
needed to meet the minimum.

\begin{center}
\begin{tabular}{@{}rrr@{}}
\toprule
Pool amount & Policyholder payment & Amount left\\
\midrule
80 & 90 & $-10$\\
100 & 90 & 10\\
120 & 108 & 12\\
\bottomrule
\end{tabular}
\end{center}

Moving from 80 to 100 leaves 20 more after payment: the ordinary share rises
by 2 and the extra floor payment falls by 18. Moving from 100 to 120 leaves
only 2 more because the minimum no longer binds. The negative amount at 80
is a funding shortfall for this component. These unweighted settlement
outcomes explain the exposure; they do not establish insurance-contract
eligibility, a present value or the insurer's profit. Other benefits and costs
remain outside this component example, and its held-pool assumption does not
replace the standard's separate accounting for investment assets
\citep[paragraphs B66(a), B107--B108]{aasb17}.

The accounting calculation therefore needs more than this payoff schedule.
Paragraph B104 defines the variable fee as the entity's share of underlying
fair value less fulfilment cash flows that do not vary based on underlying
returns. Changes in the obligation to pass through the underlying fair value
do not adjust CSM, whereas the entity-share changes relate to future service.
Financial-guarantee changes have their own future-service treatment under
B113(b). Paragraph 45 retains the loss-component and risk-mitigation exceptions,
as well as new-contract, currency and service-allocation movements. Applying
those rules requires the group's measured cash flows, discount and risk
adjustments, remaining CSM and loss component, service information and any
eligible mitigation designation. The original increase of 2 remains a valid
share-component calculation; the complete CSM movement needs these additional
inputs \citep[paragraphs 32--38, 45, 48--50 and B104--B118]{aasb17}.

The claims-development example above supplies an estimation input, not a
recognition rule. Its point projection is 148, but the single observed final
factor leaves its variance unresolved. In the general method the factors are
\[
 \widehat f_k=\frac{\sum_{i\in O_k} C_{i,k+1}}{\sum_{i\in O_k}C_{i,k}},
\]
where $O_k$ contains origin years with both required observed cells. Each
incomplete row is projected through the remaining factors. The same triangle
therefore has ultimate claims 180, 216 and 252, with unpaid amounts 0, 36 and
112. Tail development, calendar inflation, zero or negative cells and dependence
need explicit extensions. The conditional means, variance scaling and
independent-origin-year assumptions behind the inspected Mack error formulas
\citep[equations (1)--(3), Theorems 1--3]{mack1993} cannot turn this tiny
triangle into a predictive law or a justified insurance risk adjustment.

\clearpage
\section{Acquisition costs before the renewal contracts exist}
\label{sec:insurance-acquisition-recoverability}

An insurer may pay acquisition costs that help obtain both current contracts
and their expected renewals. Some of the related future groups do not yet
exist as recognised insurance balances. The retained AASB~17 provisions
therefore require a separate acquisition-asset record for each related group,
unless the permitted immediate-expense election is used. Qualifying costs
are those directly attributable to the insurance portfolio; this is not a
general permission to capitalise selling expenses. A systematic and rational
allocation distinguishes the original group from groups expected to contain
its renewals \citep[paragraphs 28A--28B, B35A and Appendix A]{aasb17}.

Consider a separate, authored example in currency units. Follow only 60 of
qualifying costs already paid and allocated to expected renewals: 40 from
one origination cohort to future group G1, and 20 from another cohort to
future group G2. Other acquisition allocations are outside this example.
Assume that both allocations are supported, neither future group has yet
been recognised, the immediate-expense election is not used, and facts now
indicate possible impairment. Those facts trigger the reporting-date
recoverability assessment under paragraph 28E. An expected renewal remains
a forecast, not a newly recognised premium receivable
\citep[paragraphs 28B, 28E and 35]{aasb17}.

Each future group contains the relevant renewals together with other new
contracts. For G1, suppose the whole group's expected net inflow is 35,
of which the renewal cohort contributes 25. For G2, the corresponding
amounts are 25 and 10. These are conditional measurements under paragraph
32(a), incorporating the required cash-flow, time-value and risk adjustments.
They are neither gross premiums nor CSM. The acquisition cash already paid
does not become a second future outflow. In this additive example, the
other contracts contribute 10 to G1 and 15 to G2.

Paragraph B35D(a) first tests the asset against its related group's expected
net inflow. G1's asset of 40 exceeds 35, so recognise impairment of 5 and
retain 35. G2's asset of 20 is below its whole-group inflow of 25, so this
test produces no loss for G2. Passing that first test does not settle whether
the renewals themselves can recover the costs allocated to them.

The second test addresses that narrower population. Under B35D(b), the
allocated renewal costs exceed the expected renewal inflows by $40-25=15$
for G1 and $20-10=10$ for G2. But 5 of G1's shortfall has already been
recognised. B35D(b)(ii) requires that overlap to be deducted, leaving
additional losses of 10 for G1 and 10 for G2. Total impairment is therefore
$5+10+10=25$, and the remaining assets are 25 and 10. Adding the first loss
to the two unadjusted renewal shortfalls would instead record 30, counting
the same loss of 5 twice \citep[paragraph B35D]{aasb17}.

For these disjoint allocations, write $A_g$ for the original asset, $N_g$ for
the whole-group expected net inflow and $U_g$ for the relevant renewal
inflow. With no prior impairment, the two losses and the resulting balance
are
\[
 I_g=\max(A_g-N_g,0),\qquad
 J_g=\max(A_g-U_g-I_g,0),\qquad
 B_g=A_g-I_g-J_g.
\]
The subtraction of $I_g$ inside the second test preserves the already
recognised loss. This local calculation relies on the identified, nonoverlapping
renewal cohorts and positive inputs. A pool with several overlapping origins
would need an explicit allocation of the relevant costs, inflows and prior
losses before the same reasoning could be applied.

\clearpage
The entries also distinguish impairment from a cash payment. Debit impairment
expense 5 and credit the G1 acquisition asset 5 for the first test. Then debit
impairment expense 20 and credit each group's asset by 10 for the second.
Cash does not move: the acquisition costs were already paid. For a small
balance-sheet slice starting with cash and equity of 100 and no liabilities,
paying the allocated costs leaves cash 40 and acquisition assets 60. After
impairment, cash remains 40, acquisition assets are 35 and equity is 75.
Paragraph 79 includes these acquisition assets within the related issued-contract
portfolio carrying amounts; the separate records used here do not prescribe
a separate line on the published balance sheet
\citep[paragraphs 28E and 79]{aasb17}.

Recoverability can improve. At a later reporting date, suppose the full-group
net inflows are 45 and 30, with renewal contributions of 35 and 15. Assume no
intervening cash movement or recognition of either group. The allocation
assumptions have been reviewed under B35B and remain unchanged; only the
measurement inputs improve. Against the original allocations of 40 and 20,
the remaining shortfalls are now 5 for each group. Paragraph 28F permits
reversal to the extent the impairment conditions have improved. Reverse 10
for G1 and 5 for G2: debit the acquisition assets 15 and credit the impairment
reversal in profit or loss 15. Assets become 35 and 15, totalling 50; cash
stays 40 and equity becomes 90. The improvement restores 15 of the earlier
loss of 25 without creating cash or recognising the expected renewals early
\citep[paragraphs 28F and B35B]{aasb17}.

If the impairment conditions later disappear altogether, with these other
facts unchanged, only the remaining loss of 10 can be reversed; the original
allocated cost of 60 is the ceiling in this example. When related contracts
are actually recognised and the acquisition cash flows enter their group
measurement, the corresponding acquisition asset must be derecognised.
B35C provides for partial incorporation when contracts are added over several
periods. Keeping both the old asset and its incorporated amount would double
count the same cost. The record must therefore preserve its origination cohort,
future group, allocation history, impairment and reversal history, and the
amount already incorporated into recognised contracts
\citep[paragraphs 28C--28F and B35B--B35C]{aasb17}.

These are ordinary reporting-date impairment and reversal entries. They do
not use the separate transition rules for acquisition assets. The historical
edition, actual cost eligibility and supported renewal estimates must be
checked before applying the example to an insurer.

\clearpage
\section{When the acquisition asset enters the insurance group}
\label{sec:insurance-acquisition-incorporation}

Continue from the preceding reversal: cash is 40, the two acquisition assets
are 35 and 15, and equity is 90. The next event is recognition of the related
insurance groups. Assume that coverage begins and the first premiums become
due at this date, with no earlier recognition trigger. All contemplated
contracts are admitted together to their correctly identified groups. No
service has yet elapsed and no claims have been incurred. Paragraphs 25--28
determine recognition; the later cash receipt is a separate event
\citep[paragraphs 25--28]{aasb17}.

Keep the measured group net inflows of 45 and 30 from the previous example.
To expose their components, suppose premiums of 100 and 60 are due immediately,
the present values of future claims and expenses are 50 and 25, and each
group has a nonfinancial risk adjustment of 5. Thus fulfilment cash flows
before premium receipt are $50+5-100=-45$ and $25+5-60=-30$, using positive
amounts for net outflows. These are supplied conditional measurements;
present values do not imply a zero discount rate. Previously paid acquisition
costs are not another future payment. Assume general measurement and no other
assets, liabilities, taxes, reinsurance or intervening measurement changes.

Paragraph 38 requires initial CSM to prevent income or expense from the FCF,
any initial cash flows and the derecognition of the acquisition asset, subject
to its stated exceptions. The paid cost therefore affects initial CSM even
though it is no longer a future cash flow. Let $A$ denote the acquisition
asset's remaining carrying amount, $F$ the FCF and $C$ the CSM. Removing the
old asset without a gain or loss requires the new group balance to replace it.
With positive $L$ denoting a liability, the relation in this case is
\[
 L=F+C=-A,\qquad C=-F-A.
\]
G1 consequently has CSM $45-35=10$ and a net group asset of 35; G2 has CSM
$30-15=15$ and a net group asset of 15. Both margins are positive, so the
onerous-group exception does not arise. There are no other previously
recognised cash-flow assets to which paragraph B123A's revenue-and-expense
exception would apply \citep[paragraphs 28C, 32, 38, 47 and B123A]{aasb17}.

An illustrative recognition entry debits the two group assets by 35 and 15
and credits the corresponding acquisition assets by those amounts. Cash is
still 40 and equity is still 90. The separate acquisition balances are now
zero. The amounts transferred are the carrying values after impairment and
reversal, rather than the original allocations of 40 and 20. The earlier
loss history remains available even though those assets have entered the
group measurements \citep[paragraphs 28C--28D and B35C]{aasb17}.

Now receive the premiums in full, immediately after that measurement. Debit
cash 160 and credit the group carrying balances by 100 and 60. Because the
premiums are no longer future inflows, FCF becomes 55 for G1 and 30 for G2.
CSM stays 10 and 15: no service or time has passed. The groups therefore carry
liabilities of $55+10=65$ and $30+15=45$. Cash has increased to 200, and
$200-65-45=90$ confirms unchanged equity. Paragraph 83 ties revenue to
service, while B123(a)(i) expressly distinguishes premium cash inflows from
the liability reductions that generate revenue. Later acquisition-recovery
revenue and the equal service expense are allocated over time under B125;
neither is created by this instantaneous receipt
\citep[paragraphs 40--41, 83--85 and B123--B125]{aasb17}.

\clearpage
The three balance-sheet states make the transfer visible. Amounts are
currency units; group assets and liabilities below are their net carrying
amounts, not their separate FCF and CSM components.
\begin{center}
\begin{tabular}{lrrr}
\toprule
 & Before recognition & After recognition & After receipt \\
\midrule
Cash & 40 & 40 & 200 \\
Acquisition assets & 50 & 0 & 0 \\
Group assets & 0 & 50 & 0 \\
Group liabilities & 0 & 0 & 110 \\
Equity & 90 & 90 & 90 \\
\bottomrule
\end{tabular}
\end{center}
Receiving 160 has changed liquidity and the insurance balances without
creating profit. Omitting the acquisition assets from the initial CSM
calculation would instead give margins of 45 and 30, overstating total CSM
by 50. Retaining the separate acquisition assets after the correct transfer
would overstate assets and equity by 50.

Even a balanced statement can conceal both mistakes together. If the
implementation leaves acquisition assets of 50 and uses the overstated
margins, its post-receipt liability becomes 160. Cash 200 plus the retained
asset 50 less liability 160 still equals equity 90. That identity passes,
but the asset, liability and CSM are wrong relative to the stated measurement.
Checking the components and the derecognition event is therefore necessary
in addition to checking that the statement balances.

As a boundary case, if no acquisition asset exists, $A=0$ gives $C=-F$ and
a zero group carrying amount before premium receipt, under the same
positive-margin assumptions. A general implementation must also handle
onerous groups, other initial cash-flow balances and contracts added over
time. The present calculation establishes the two-event bridge for the
specified groups, rather than those additional cases or an insurer's actual
valuation.

\clearpage
\section{Reading acquisition costs in an actual report}
\label{sec:munich-new-contracts}

The preceding example supplied the acquisition assets and the net inflows
for two identified groups. A published insurer's accounts reveal a different
level of detail. Munich Re's 2023 report shows contracts recognised for the
first time during the year, separating concluded from acquired contracts and
nononerous from onerous business. Note 27 covers segments predominantly using
general measurement or the variable fee approach; it omits two predominantly
PAA segments. These are segment disclosures, not a list of individual groups
\citep[printed pp. 252--254]{munichre2023annual}.

Consider the ERGO Life and Health Germany segment, identified in the paragraph
preceding its table. Its newly recognised contracts combine GMM and VFA
measurements. The following two columns reproduce its concluded-contract
amounts in millions of euros. Preserve the report's signs: inflows are positive,
outflows and the margin are negative. Its onerous CSM cell is blank; a dash
below preserves that blank instead of assigning a reported zero
\citep[printed pp. 252--253]{munichre2023annual}.

\begin{center}
\begin{tabular}{p{0.58\textwidth}rr}
\toprule
 & Nononerous & Onerous \\
\midrule
Present value of future inflows & 2,842 & 17 \\
Future claims, expenses and investment components & $-2,438$ & $-12$ \\
Future acquisition costs & $-243$ & $-25$ \\
Risk adjustment & $-8$ & 0 \\
Contractual service margin & $-154$ & --- \\
Reported total & 0 & $-21$ \\
\bottomrule
\end{tabular}
\end{center}

The nononerous amounts illustrate how expected future profit is deferred in
CSM. The onerous column retains a loss-bearing population separately. Munich
Re describes recognising an expected initial loss in profit or loss and
keeping onerous contracts separate from profitable ones. Combining these
columns into one net profit would erase that distinction
\citep[printed pp. 165--166]{munichre2023annual}.

Adding the displayed nononerous numbers gives $-1$, although the table reports
zero. The onerous cash-flow and risk-adjustment rows give $17-12-25-0=-20$,
against a reported total of $-21$. The report states that amounts are rounded
to millions and that totals may consequently differ. A possible set of
unreported nononerous values is $2{,}842.3-2{,}437.8-242.8-7.8-153.9=0$;
each rounds to its displayed amount. Similarly,
$16.6-12.1-25.2-0.3=-21$ is consistent with the onerous column if no CSM is
assigned. Together these hypothetical values also reproduce the table's
outflow subtotals and aggregate rows after rounding. They demonstrate
consistency under nearest-million rounding; they do not recover the issuer's
unreported measurements or establish its exact rounding convention. Retain
the original figures and blank cells \citep[printed pp. 155, 253]{munichre2023annual}.

\clearpage
The acquisition-cost row appears within \emph{future cash outflows}. Its
reported total of $243+25=268$ therefore cannot simply be entered as the
previous example's asset $A$, which arose from costs already paid before the
related groups were recognised. Nor do the table's acquired-contract columns
describe how sales acquisition costs were paid. Those columns concern a
different distinction between contract populations. Their displayed zeros
retain the same rounding qualification. The selected disclosure does not
identify separate prior acquisition assets, their impairment history or their
allocation to individual groups \citep[printed pp. 252--253]{munichre2023annual}.

A third quantity appears in the segment's service revenue. For 2023, Munich
Re reports €202 million as the portion of premium relating to recovery
of acquisition costs. Its policy explains that acquisition costs enter
insurance service expense in the same amount as the recovery included in
revenue, allocated over the coverage period. This is a period service amount;
the €268 million above is an estimate of future acquisition outflows for
newly recognised contracts. Subtracting one from the other would not produce
a remaining acquisition asset. The relevant populations, dates and movements
must first be linked \citep[printed pp. 177, 253, 264]{munichre2023annual}.

The segment's liability reconciliation makes that caution concrete: it shows
€637 million in the service-expense row labelled ``Recovery of acquisition
costs'', and €978 million of acquisition costs within cash flows. Thus the
€202 million recovery in revenue does not describe the entire expense row.
The difference of €435 million is an arithmetic residual; assigning its
cause requires the relevant populations and expense treatments. Neither that
residual nor cash less expense identifies a separate acquisition asset
\citep[printed pp. 230--231]{munichre2023annual}.

The same timing distinction applies to profit. The report gives
€819 million of CSM for services provided during 2023 in this segment.
That amount concerns service during the year, whereas €154 million is
the margin attached to the newly recognised contracts in the selected table.
It is not a release from that new-business amount alone. A reconciliation
needs the opening CSM and every intervening movement. The segment's total
insurance revenue also combines measurement approaches: €7,966 million
under GMM/VFA and €1,976 million under PAA make the reported
€9,942 million \citep[printed pp. 253, 264]{munichre2023annual}.

Even within one insurer, an acquisition-cost policy needs its applicability
conditions. Munich Re reports using the PAA option to expense acquisition
costs directly for a substantial part of its property-casualty primary
insurance business. The earlier general-measurement asset example therefore
requires the relevant treatment and group facts before it can be applied to
a real portfolio \citep[printed p. 166]{munichre2023annual}. These are readings
of the retained 2023 English translation, whose German original is binding.
They identify reported aggregates and necessary further facts; the underlying
group histories and valuations remain unresolved.

\clearpage
\section{Adoption dates and the restatement of 2022}
\label{sec:insurance-adoption}

A standard's effective date and the economic date of a reported balance need
not coincide. An insurer adopting new rules in 2023 may report a revised
balance for December 2022 while retaining the earlier report as a record of
what it originally disclosed. The system must preserve both versions and
explain which one a calculation uses.

The Ministry of Finance's explanation dated 12 June 2025 describes the original
ASBE~25 timetable. It specifies application from 1 January 2023 for enterprises
listed both domestically and overseas, and for overseas-listed enterprises
preparing financial statements under IFRS or ASBE. Other enterprises applying
ASBE were scheduled to apply the revised standard from 1 January 2026.
Its account of the
2025 notice says existing adopters, including enterprises that had adopted
early under the 2020 notice, must continue applying it. Insurers facing genuine
difficulties and seeking deferral must submit written reasons to the Ministry's
Accounting Department and the relevant institutional supervision departments
of the National Financial Regulatory Administration. That statement does not
establish an automatic deferral merely because preparation is difficult
\citep[answers on background and implementation arrangements]{mofinsurance2025qa}.

The same explanation limits the new simplification choices to unlisted
enterprises; matters outside those choices remain subject to ASBE~25. It names
grouping, measurement, presentation and transition as affected areas, including
comparative information and the transition date. These are source-supported
descriptions of the Ministry's explanation. The operative 2020 and 2025 notices
have not yet been directly verified in this collection, so the explanation
does not supply the detailed conditions needed to implement those choices
\citep[answer on the relationship between the notice and the standard]{mofinsurance2025qa}.

Ping An provides a concrete adoption history. Its 2022 ASBE annual report
states that the group would adopt the revised insurance standard for accounting
years beginning on 1 January 2023. Its 2023 report confirms adoption on that
date and restatement of comparative information
\citep[Note XIV.2, pp.~315--316]{pingan2022asbe}
\citep[Note IV.34, p.~192]{pingan2023asbe}. The following balances all refer
to the consolidated group at 31 December 2022, in RMB millions. The first
column comes from the 2022 report; the second is the restated comparative in
the 2023 report. Changes are calculated from those columns; the asset,
liability and parent-equity changes also match the three rows of Note IV.34's
reported bridge \citep[pp.~146--148]{pingan2022asbe}
\citep[pp.~141--143 and 192]{pingan2023asbe}.

\begin{center}
\small
\begin{tabularx}{\linewidth}{@{}Xrrr@{}}
\toprule
Balance & Original & Restated & Change\\
\midrule
Assets & 11,137,168 & 11,009,940 & $-127,228$\\
Liabilities & 9,961,870 & 9,823,944 & $-137,926$\\
Total equity & 1,175,298 & 1,185,996 & $10,698$\\
Parent shareholders' equity & 858,675 & 869,191 & $10,516$\\
Non-controlling interests & 316,623 & 316,805 & $182$\\
\bottomrule
\end{tabularx}
\end{center}

The distinction between total equity and equity attributable to the parent
matters. Assets less liabilities rise by
$-127,228-(-137,926)=10,698$. Of that increase, 10,516 belongs to parent
shareholders and 182 to non-controlling interests. Treating the parent-equity
row as total equity would leave an unexplained difference of 182. These are
changes between reported accounting bases for the same date, not cash received
in a later period.

The transition also changes how later growth is calculated. Ping An reports
assets of 11,583,417 at the end of 2023. Subtracting its restated 2022
comparative gives a reported-balance increase of
$11,583,417-11,009,940=573,477$. Subtracting the original 2022 amount instead
gives 446,249, mixing the old and new bases and understating that comparable
movement by 127,228 \citep[p.~141]{pingan2023asbe}. The first difference is
a movement in reported asset balances; it is not a cash-flow measure or an
explanation of every economic cause of growth.

An original-source query for assets at 31 December 2022, explicitly selecting
the 2022 report, therefore returns 11,137,168. A query selecting the 2023
report's restated comparison for that date returns 11,009,940. Each assertion
retains its report identity, consolidated scope, RMB-million scale, accounting
basis and economic date. A link records that the later assertion restates the
comparison; it does not delete the original. A request for a comparable annual
movement must reject the mixed-basis pair unless an explicit reconciliation
has first aligned it. Because the basis and disclosure version differ, these
two amounts are not a same-context source conflict.

The information date remains a separate question. The 2023 report says its
board authorised the statements on 21 March 2024
\citep[Note I, p.~154]{pingan2023asbe}. Neither that authorisation nor the
31 December reporting date establishes when the exact retained document became
public. The corpus has archive-date leads but no verified publication date or
historical byte vintage for either selected report. A query for information
available to investors at 31 December 2023 must therefore retain an unresolved
availability result for this source selection; it cannot silently import the
later restated balance. The corpus's October 2026 revalidation timestamps
describe local evidence handling, not investor knowledge in 2023.

Finally, Note IV.34 reports financial-asset reclassification, including use of
the classification overlay for relevant assets derecognised during 2022. It
also states that some insurance groups used modified-retrospective or
fair-value transition because retrospective application was impracticable.
The headline bridge consequently cannot be reconstructed by applying one
percentage to every insurance liability. The example establishes how to retain
and compare disclosed histories; reconstructing the transition itself still
requires the affected groups, asset classifications and underlying measurement
inputs \citep[Note IV.34, p.~192]{pingan2023asbe}.

\clearpage
\section{Reconstructing balances at transition}
\label{sec:insurance-transition}

A disclosed restatement tells us what balances an insurer reported under the
new basis. It does not reveal the historical contract groups, assumptions or
services needed to reproduce their opening measurement. The three transition
methods address different limits on that reconstruction. The retained
December 2022 AASB~17 compilation provides a useful, explicitly scoped account
of those methods: its stated applicability covers annual periods beginning
from 1 January 2023 to before 1 July 2026. These provisions explain a design
requirement; they do not establish the rules governing Ping An's ASBE groups
\citep[Appendix C]{aasb17}.

First distinguish the dates. Initial application begins the annual period in
which the entity first applies the standard. Transition normally begins the
immediately preceding annual period. Thus a calendar-year entity first applying
the standard in 2023, with only 2022 presented as an adjusted comparative,
starts its transition calculation on 1 January 2022. If it elects to present
earlier adjusted comparatives, paragraph C25 instead takes transition back to
the beginning of the earliest such period. Neither date establishes when the
resulting report became public \citep[paragraphs C2 and C25--C27]{aasb17}.

Full retrospective application identifies and measures each group as if the
standard had always applied. Acquisition-cash-flow assets are considered
separately, balances that would not exist under the new rules are removed, and
the resulting net difference is recognised in equity. This route needs the
historical information required to reconstruct grouping, cash flows,
measurement assumptions and service already provided. It is required unless
application is impracticable or the specific paragraph C5A exception applies
\citep[paragraphs C3--C5]{aasb17}.

Impracticability requires more than a missing historical workbook. Under
AASB~108, the entity must have made every reasonable effort to apply the
requirement. Relevant difficulties include effects that cannot be determined,
assumptions that would have to be made about management's past intentions,
and significant estimates for which the necessary contemporaneous evidence
cannot be objectively distinguished from other information. Those conditions
need to be established for the reconstruction at issue
\citep[paragraph~5]{aasb108comp2022}.

An estimate prepared later can nevertheless describe an earlier period. The
question is whether it reflects the circumstances that existed then, using
evidence that would have been available when the earlier financial statements
were authorised for issue. A newly assembled estimate can meet that test;
replacing an earlier expectation with knowledge of what subsequently happened
uses hindsight. The date of preparation alone therefore does not settle the
question. Nor does authorisation establish when an outside investor could
obtain the information: that requires the separate public-availability record.
The insurance standard's specific transitional provisions govern their stated
cases, including permitted information under the modified and fair-value
routes described below
\citep[paragraphs~7, 19(a) and 50--53]{aasb108comp2022}.

When full retrospection is impracticable for a group, the alternatives are
modified retrospection and fair-value transition. Modified retrospection aims
for the closest retrospective result possible using reasonable and supportable
information available without undue cost or effort. Each specified modification
is allowed only to the extent that the necessary retrospective information is
missing. If the information needed even for this modified method cannot be
obtained, paragraph C6(a) requires the fair-value approach. Missing history
therefore does not license an unrestricted choice of present assumptions
\citep[paragraphs C5--C8]{aasb17}.

For example, the modified route can reconstruct earlier expected cash flows
from estimates at transition, or an earlier determinable date, adjusted for
known intervening cash flows, including those of contracts that have expired.
Its discount-rate approximations require the specified yield-curve evidence;
its risk adjustment and allocation of past service require their own records.
The source also prescribes nil amounts for particular acquisition-cash-flow
items when the necessary information is unavailable. That narrow instruction
cannot be generalized into a rule that every missing amount is zero
\citep[paragraphs C12--C16]{aasb17}.

Fair-value transition takes a different starting point. Paragraph C20 determines
the CSM or loss component from the difference between the fair value of the
insurance group and its fulfilment cash flows, both measured at transition.
The insurance group's fair value is distinct from the value of its underlying
investment pool. Replacing one with the other would compute a different
quantity. The calculation needs a supported group valuation as well as the
cash-flow, discount and risk-adjustment inputs; C20 expressly excludes the
demand-feature rule in AASB~13 paragraph~47. Its permitted information-date,
grouping and accumulated-OCI choices also have specified conditions
\citep[paragraphs C20--C24B]{aasb17}.

To interpret that group valuation, consider what would be paid to transfer
the liability in an orderly transaction between market participants at the
measurement date. The obligation remains outstanding, and the transferee must
fulfil it; the thought experiment is not an immediate settlement with the
policyholder. AASB~13 therefore calls for assumptions that market participants
would use, in the accessible principal market or, in its absence, the most
advantageous market. The unit being valued still comes from C20: the insurance
group \citep[paragraphs~2--3, 9, 13--24 and 34]{aasb13comp2021}.

A quoted price is not required in every case. The valuation should maximise
relevant observable inputs, while supported unobservable inputs can supply
the participant assumptions for which market observations are unavailable.
The insurer's own data can be a starting point, but require adjustment when
reasonably available evidence indicates different participant assumptions or
entity-specific features unavailable to other participants. A technique must
also fit the circumstances and available data. These requirements make the
source, date and interpretation of each valuation input consequential
\citep[paragraphs~61--67 and 87--89]{aasb13comp2021}.

There is a further reason to retain separate fair-value and FCF calculations.
AASB~13 includes non-performance risk, including the issuer's own credit
risk, in liability fair value and assumes that risk is unchanged by the
hypothetical transfer. AASB~17 expressly excludes the issuer's own
non-performance risk from fulfilment cash flows for issued insurance contracts.
Even when the two measurements draw on some of the same cash-flow records,
their treatment of this risk differs. Neither a single imported FCF balance
nor the value of the underlying investments supplies both C20 inputs
\citep[paragraphs~42--44]{aasb13comp2021}
\citep[paragraph~31]{aasb17}.\footnote{The comparison uses AASB~108's
31 December 2022 compilation, applying to periods beginning from 1 January
2023, and AASB~13's 30 June 2021 compilation, applying from 1 July 2021.
These are retained historical editions; later amendments and the actual
entity's applicable requirements need their own check.}

\clearpage
There is one important qualification to the ordinary method choice. Under
C5A, a group with direct participation features may elect fair-value transition
even when retrospective application is possible, if the entity chooses the
B115 risk-mitigation option prospectively from transition and used qualifying
mitigation before that date. The applicable instrument, qualification and
designation conditions still matter. Treating every fair-value election as
proof of impracticability would therefore be wrong
\citep[paragraphs C3(b), C5A and B115--B118]{aasb17}.

The reconstruction must consequently retain the method and its basis for each
group, each modification used, the information date and evidence supporting
it, the separately measured FCF, CSM or loss component, relevant OCI and
acquisition balances, and the opening-equity reconciliation. Some transition
rules expressly permit transition-date information for specified assessments.
That permission concerns the accounting restatement; it does not make a later
report admissible in a forecast made before its publication. Detailed
impracticability judgments under AASB~108 and valuation under AASB~13, together
with actual group facts, remain necessary before this method description can
produce a supported opening measurement.


\clearpage
\section{From transition measurements to opening equity}
\label{sec:insurance-transition-balances}

The fair-value approach produces a margin or a loss component, but those two
outcomes enter the liability differently. Consider two distinct groups of
issued contracts for which the grouping and eligibility for this transition
method have already been established. Use the historical provisions just
described and assume that both groups contain only remaining coverage. All
amounts below are authored currency units at the same transition date. There
are no incurred claims, reinsurance, acquisition balances, tax, accumulated
OCI or other transition adjustments in this example. These simplifications
isolate the effect of the two measurements; they do not supply actual insurer
valuations or evidence of eligibility.

For group A, suppose the measured fulfilment cash flows are 90 and the
separately measured fair value is 110. The FCF amount already includes the
discounting and nonfinancial risk adjustment required by the measurement
rules. Paragraph C20 makes the difference, $110-90=20$, the CSM. This is
unearned profit associated with future services. Keeping it within the
liability gives $90+20=110$; recording it as immediate profit would contradict
that role \citep[paragraphs 32--38, 40(a) and C20]{aasb17}.

Now reverse the ordering of the inputs. For group B, let FCF be 120 and fair
value 100. The difference identifies a loss component of 20, while CSM is
zero. The liability is 120. Paragraphs 49--50 describe the loss component as
part of the liability for remaining coverage, with its own subsequent
allocation and reversal rules. Adding another 20 to the FCF balance would
count that component twice. Replacing the liability by fair value, 100,
would instead understate the required FCF balance
\citep[paragraphs 40, 47, 49--50 and C20]{aasb17}.

Write $F$ for measured future-service FCF, $V$ for the group's fair value,
$C$ for CSM and $K$ for the nonnegative loss component. Positive amounts
denote liabilities. The two sign cases give the following local derivation:
\[
 C=\max(V-F,0),\qquad K=\max(F-V,0),\qquad
 L=F+C=\max(F,V).
\]
Here $L$ is the recognised remaining-coverage liability. The loss component
$K$ is tracked within $L$ and is not added to it. Thus $C-K=V-F$, and at
$V=F$ both components are zero. These relations apply to the stated
transition case; later changes in FCF, service provision and finance effects
follow the subsequent-measurement rules.

The numerical comparison makes the distinction visible:
\begin{center}
\begin{tabular}{lrr}
\toprule
Currency units at transition & Group A & Group B \\
\midrule
Fulfilment cash flows & 90 & 120 \\
Fair value of the group & 110 & 100 \\
CSM & 20 & 0 \\
Loss component within the liability & 0 & 20 \\
Recognised liability & 110 & 120 \\
\bottomrule
\end{tabular}
\end{center}
The recognised liability equals fair value for A and exceeds it for B.
Calling the method fair-value transition therefore does not mean that every
resulting liability is carried at fair value.

\clearpage
To see the opening-equity effect, suppose the old carrying liabilities are
105 for A and 115 for B. Assets of 500 remain unchanged, and these are the
only liabilities. Before transition, equity is $500-105-115=280$. Afterwards,
it is $500-110-120=270$. Each group's liability increases by 5, so total
opening equity falls by 10. Paragraph C5 replaces the retrospective
measurement requirement in C4(a); the removal of incompatible old balances
and the net adjustment to equity in C4(b)--(c) still apply. The ordinary
initial-recognition loss rule does not turn this transition adjustment into
current-period operating expense \citep[paragraphs C4--C5]{aasb17}.

Illustrative component entries show the same reconciliation. For A, debit
the old liability 105 and opening equity 5, then credit the FCF liability 90
and the CSM liability 20. For B, debit the old liability 115 and opening
equity 5, then credit the FCF liability 120. Record B's loss component of 20
within that 120 for its subsequent allocation; it produces no additional
credit. Neither entry changes cash. An actual transition would also need the
acquisition, reinsurance, tax and OCI effects excluded here.

The order of calculation matters as much as the signs. If an implementation
first combined the two already distinct groups, it would find total FCF of
210 and total fair value of 210. Applying the sign rule to those totals would
produce neither CSM nor a loss component and would report a liability of 210.
The correct group-by-group calculation gives 230. Pooling would therefore
understate liability and overstate equity by 20. The programme must retain
the admitted group boundary while calculating CSM and loss components, then
aggregate the resulting balances. Whether different grouping is itself
permitted requires the separate grouping rules; an aggregation shortcut
cannot decide that accounting question.

The example supplies the arithmetic once the group's two measurements and
old balances are known. A real transition disclosure may reveal only some
of those quantities. The later Munich Re case shows both what can be read from an
insurer's report and what still requires contract and valuation records.

\clearpage
\section{The recoverable part of a transition loss}
\label{sec:reinsurance-transition-recovery}

Group B's loss component of 20 describes the issued contracts. If the insurer
also holds reinsurance, how much of that loss is expected to be recovered?
Extend the preceding example by adding a separately recognised, qualifying
group of reinsurance contracts held. Assume it covers all contracts in B and
only B, with remaining coverage and no incurred claims, and uses general
measurement. This changes the earlier no-reinsurance assumption. The issued
obligation to policyholders remains 120, with its loss component of 20;
purchasing reinsurance does not extinguish it
\citep[paragraphs 60--64 and 75]{aasb17}.

The transition calculation has a specific timing exception. Ordinarily,
paragraph B119C permits the recovery-income treatment in paragraph 66A only
if the reinsurance contract was entered into before or when the onerous
underlying contracts were recognised. Suppose the underlying contracts'
recognition date, applying these rules, is 1 July 2020, the treaty was entered
into on 1 October 2020 and
both coverages remain in force at transition on 1 January 2022. The treaty
does not meet that ordinary condition for the original underlying recognition.
Paragraph C20A nevertheless expressly removes B119C's condition when
determining the loss-recovery component at transition. The exception does
not remove the treaty's other recognition, grouping, boundary or measurement
requirements \citep[paragraphs 61--66B, B119C and C20A]{aasb17}.

Let $q$ be the percentage of claims on the underlying contracts that the
insurer expects to recover from this reinsurance group, expressed as a
fraction. Assume $q=0.60$ for this conditional illustration. A nominal treaty
quota alone does not establish that expectation: coverage, limits, exclusions
and collectability matter. The full reinsurance measurement also includes
issuer non-performance risk and the nonfinancial risk transferred, under
paragraphs 63--64. With $K=20$ still denoting B's loss component, C20A gives
the reinsurance loss-recovery component $R$:
\[
 R=Kq=20\times 0.60=12.
\]
This 12 belongs to the reinsurance asset for remaining coverage. It is not
the issued liability, the amount of expected claims, or a valuation of the
whole reinsurance asset. If only some contracts in the underlying group were
covered, the required allocation of the losses would also have to be
established; the all-covered assumption avoids that extra calculation here
\citep[paragraphs 63--64, 66B, B119E and C20A]{aasb17}.

The complete asset still needs the treaty cash flows, discounting, transferred
risk and net cost or gain of purchasing reinsurance. The reinsurance CSM
reflects that net cost or gain; held reinsurance is not classified as onerous
under paragraph 68. Applying the issued-group CSM and loss formulas to the
held contract would therefore be wrong. Presentation likewise keeps the
issued and held-reinsurance portfolios separate
\citep[paragraphs 65--68 and 78--79]{aasb17}.

\clearpage
To extend the opening balance sheet, supply the missing total explicitly.
Assume a separate complete measurement gives a reinsurance asset of 50,
already including the loss-recovery component of 12, against an old asset of
45. These are additional conditional inputs; the multiplication above does
not derive them. Other assets remain 500 and the two issued liabilities still
increase from 220 to 230. Total assets therefore rise from $500+45=545$ to
$500+50=550$, while equity moves from $545-220=325$ to $550-230=320$.
The asset increase adds 5 to equity against the issued-liability increase
of 10, giving a net reduction of 5. An illustrative reinsurance entry debits
the new asset 50, credits the old asset 45 and credits opening equity 5.
Adding another asset of 12 would count a component of the supplied 50 twice
and overstate equity by 12. The net transition change belongs in equity
under C4(b)--(c), rather than being inferred as current recovery income
from the ordinary recognition rule \citep[paragraphs C4--C5 and C20A]{aasb17}.

The component also imposes a useful later check. If B's loss component falls
to 14 and the recoverable portion remains 60\%, paragraph B119F limits the
reinsurance loss-recovery component to $14\times0.60=8.4$. Leaving it at 12
would violate that bound. The bound alone neither selects a unique new
component nor determines the full reinsurance asset or period's profit;
the actual changes and allocation rules still matter. Similarly, $q=0$
gives no loss-recovery component and $q=1$ gives 20, without erasing the
issued obligation or determining the entire asset
\citep[paragraphs 66--67 and B119F]{aasb17}.

\clearpage
\hypertarget{munich-re-transition-disclosure}{%
\section{What an insurer's transition disclosure reveals}\label{method-insurance-munich-transition}}

Munich Re's 2023 report shows how much a public disclosure can resolve.
Its consolidated statements apply IFRS as adopted by the European Union,
together with the stated German commercial-law requirements. The first
application of IFRS~17 is 1 January 2023, with transition at 1 January 2022
and restated insurance comparatives. Financial-instrument comparatives
generally remain under IAS~39, subject to the disclosed classification
overlay. A single year label therefore does not imply a single measurement
basis for every line in the comparative balance sheet
\citep[pp.~155 and 183--184]{munichre2023annual}.\footnote{The retained
English report states that it is an informational translation and the German
original is binding (p.~1). The discussion attributes the disclosed policies
and amounts; the historical AASB texts above are comparison sources.}

The report identifies full retrospection for certain property-casualty
contracts and parts of life reinsurance. Where historical information is
insufficient in life reinsurance, it mostly uses fair-value transition; it
also uses that method for specified property-casualty reinsurance groups and
German primary business. Modified retrospection and fair value together
account for over 90\% of net insurance liabilities at transition. That is a
statement about a liability denominator at a particular date. It supplies
neither a count of contracts nor the assignment and supporting history of
every group \citep[p.~183]{munichre2023annual}.

The opening-equity bridge makes the scale of the different effects visible.
The reported margin component is much larger than the net equity reduction
because the other measurement changes also enter the reconciliation. The
following amounts reproduce the report's signs and billion-euro scale
\citep[p.~184]{munichre2023annual}.

\begin{center}
\begin{tabular}{lr}
\toprule
Opening-equity reconciliation at 1 January 2022 & EUR billions \\
\midrule
Equity under IFRS 4 & 30.9 \\
Modified measurement basis & 25.3 \\
Contractual service margin & $-22.3$ \\
Risk adjustment for non-financial risk & $-5.6$ \\
Loss component & $-1.4$ \\
Deferred taxes & 1.5 \\
Reported change in equity & $-2.5$ \\
Equity under IFRS 17 & 28.5 \\
\bottomrule
\end{tabular}
\end{center}

The five displayed change components sum to $-2.5$ billion, whereas the
displayed endpoints differ by $-2.4$ billion. The report warns that rounding
can cause such deviations; these figures are compatible under rounding to
the displayed precision. The representation should retain both the reported
endpoints and the bridge, with their precision, instead of changing one to
force an exact identity \citep[pp.~1, 155 and 184]{munichre2023annual}.

\clearpage
The separate margin tables also clarify which contracts are being followed.
At 1 January 2022 the reported issued-contract CSM is $-22{,}688$ million
and the held-reinsurance net cost/gain is $368$ million. Adding those reported
totals gives $-22{,}320$ million, consistent with the rounded $-22.3$ billion
margin component above. The issued-contract transition table follows groups
in force at transition, and its CSM amounts concern the general measurement
model and VFA; the associated revenue also includes PAA contracts. Its
later balances therefore cannot be substituted for the whole subsequent
portfolio or used as if CSM and revenue had identical populations
\citep[pp.~210--211 and 254--255]{munichre2023annual}.

Finally, disclosure of a method does not supply all its inputs. Munich Re
reports a 6\% cost-of-capital rate for its non-financial risk adjustment,
applied to projected capital requirements with diversification and discounting.
Its equivalent one-year confidence level is about 90\% at transition, compared
with 88\% for 2023 and 87\% for the comparative period. These are dated
properties of the insurer's calculation, not universal IFRS parameters.
The rate also cannot be imported into the separate fair-value calculation
merely because that discussion mentions cost of capital. Reconstructing an
individual group still requires its cash flows, capital projections, curves,
allocations and supported transition-method choice
\citep[pp.~165, 183 and 185]{munichre2023annual}.

\clearpage
\hypertarget{interfaces-exclusions-and-acceptance-cases}{%
\section{Interfaces, exclusions and acceptance cases}\label{method-insurance-methods-interfaces-exclusions-and-acceptance-cases}}

Insurance events are group\_recognize, premium\_settle, service\_release, claim\_incur, claim\_settle, assumption\_change, onerous\_change and treaty\_change, each bound through the event-basis, authorization and commit contracts to the exact group/version, evidence and information view. Journals and component roll-forwards must reconcile independently; a component reallocation is not always a net liability change. Invalid grouping, missing rate/law, unsupported option valuation, ambiguous contract boundary or unknown service units produces the corresponding unresolved/outside-scope result. Failed numerical paths retain their coverage/mass, without successful-only normalization. Modified contracts need the applicable derecognition analysis; no new ID can erase unfavorable old-group obligations.

Initial transition requires an actual method choice and eligibility for full retrospective, modified-retrospective or fair-value transition under AASB~17 Appendix C, with acquisition, grouping, past service and historical-rate records. This collection does not claim an executed transition or a general VFA valuation. The nonzero-rate general case and onerous allocation through expiry above are finite specification examples. Required later acceptance still includes a priced guarantee, acquisition/reinsurance interaction, transition and actuarial backtesting. These retained obligations prevent the two-year illustration from being described as full insurance-domain readiness.
